\chapter{Installation}
\label{chap:installation}

\section{Linux}

\subsection{Paquet Debian / Ubuntu (\texttt{.deb})}

Téléchargez le fichier \texttt{pomone\_<version>\_amd64.deb} depuis la page \href{https://github.com/guycorbaz/pomone/releases}{Releases}, puis installez-le :

\begin{lstlisting}[language=bash]
sudo dpkg -i pomone_0.1.0_amd64.deb
sudo apt-get install -f   # complete dependencies if needed
\end{lstlisting}

Le binaire est installé dans \texttt{/usr/bin/pomone}, avec une entrée \texttt{.desktop} et les icônes système. Pomone apparaît dans le menu d'applications de votre environnement de bureau.

Dépendances système requises (installées automatiquement par \texttt{apt}) :

\begin{itemize}
  \item \texttt{libfontconfig1} (ou \texttt{libfontconfig1t64} sur Ubuntu 24.04+)
  \item \texttt{libfreetype6}
  \item \texttt{libxkbcommon0}
\end{itemize}

\subsection{AppImage}

L'AppImage est un format autonome qui ne nécessite aucune installation. Téléchargez \texttt{pomone\_<version>\_x86\_64.AppImage}, rendez-le exécutable et lancez-le :

\begin{lstlisting}[language=bash]
chmod +x pomone_0.1.0_x86_64.AppImage
./pomone_0.1.0_x86_64.AppImage
\end{lstlisting}

\subsection{Compilation depuis les sources}

\begin{lstlisting}[language=bash]
git clone https://github.com/guycorbaz/pomone
cd pomone
sudo apt install libfontconfig-dev pkg-config   # Debian/Ubuntu
cargo build --release
./target/release/pomone
\end{lstlisting}

Prérequis : \texttt{rustc} 1.80 ou plus récent.

\section{Windows}

\begin{warnbox}
Le packaging Windows (\texttt{.msi}) est prévu pour une version ultérieure. En attendant, vous pouvez compiler depuis les sources avec \href{https://rustup.rs/}{rustup} installé.
\end{warnbox}

\section{macOS}

\begin{warnbox}
Le packaging macOS (\texttt{.dmg}) est prévu pour une version ultérieure. En attendant, vous pouvez compiler depuis les sources avec \href{https://rustup.rs/}{rustup} installé.
\end{warnbox}

\section{Emplacement des données}

Par défaut, Pomone crée sa base de données SQLite au premier lancement à l'emplacement suivant :

\begin{itemize}
  \item \textbf{Linux} : \texttt{\~/.local/share/pomone/pomone.sqlite}
  \item \textbf{Windows} : \texttt{\%APPDATA\%\textbackslash pomone\textbackslash pomone.sqlite}
  \item \textbf{macOS} : \texttt{\~/Library/Application Support/pomone/pomone.sqlite}
\end{itemize}

Vous pouvez changer l'emplacement ou basculer vers un serveur MariaDB depuis l'écran \emph{Paramètres} (voir chapitre \ref{chap:parametres}).
